% Part: intuitionistic-logic
% Chapter: tableaux
% Section: introduction

\documentclass[../../../include/open-logic-section]{subfiles}

\begin{document}

\olfileid[hi]{int}{tab}{int}

\olsection{परिचय}

!!{tableau}, !!{signed formula}ों के कुछ (नीचे की ओर शाखित) वृक्ष हैं;
अर्थात ऐसे युग्म जिनमें सत्य-मूल्य चिह्न ($\True$ या $\False$) और कोई
!!{sentence} होता है:
\[
\sFmla{\True}{!A} \text{ अथवा } \sFmla{\False}{!A}.
\]
कोई !!{tableau} कुछ \emph{मान्यताओं} से आरंभ होता है। प्रत्येक अगला
!!{signed formula} किसी अनुमान-नियम को लगाकर उत्पन्न किया जाता है। कुछ
अनुमान-नियम वृक्ष के किसी सिरे पर एक या अधिक !!{signed formula} जोड़ते हैं;
अन्य दो नए सिरे जोड़कर दो शाखाएँ बनाते हैं। नियम ऐसे !!{signed formula} देते
हैं जिनका !!{formula}, उस !!{signed formula} के सूत्र से कम जटिल होता है जिस
पर नियम लगाया गया था। जब किसी शाखा में $\sFmla{\True}{!A}$ और
$\sFmla{\False}{!A}$ दोनों हों, तो हम शाखा को \emph{संवृत} कहते हैं। यदि
किसी !!{tableau} की प्रत्येक शाखा संवृत हो, तो पूरा !!{tableau} संवृत है। कोई
संवृत !!{tableau} ऐसी !!{derivation} बनाता है जो दिखाती है कि टैब्लो आरंभ
करने के लिए प्रयुक्त !!{signed formula}ों का समुच्चय असंतोषणीय है। इससे
$\Proves$ संबंध परिभाषित किया जा सकता है: $\Gamma \Proves !A$ तभी जब कोई
परिमित समुच्चय~$\Gamma_0 = \{!B_1, \dots, !B_n\} \subseteq \Gamma$ ऐसा हो
कि निम्न मान्यताओं के लिए कोई संवृत !!{tableau} हो:
\[
\{\sFmla{\False}{!A}, \sFmla{\True}{!B_1}, \dots, \sFmla{\True}{!B_n}\}.
\]

अंतःप्रज्ञावादी तर्कशास्त्र के लिए हमें !!{signed formula} की धारणा को
विस्तृत करना और संयोजकों के नियमों को समायोजित करना होगा। किसी चिह्न
($\True$ या $\False$) के अतिरिक्त, मोडल !!{tableau} के !!{formula}ों में
\emph{उपसर्ग}~$\sigma$ भी होते हैं। उपसर्ग धन पूर्णांकों के अरिक्त अनुक्रम
हैं; अर्थात $\sigma \in (\PosInt)^* \setminus \{\emptyseq\}$। ऐसे उपसर्गों
को हम आसपास के $\tuple{\ }$ के बिना लिखते हैं और अलग-अलग !!{element}ों को
अल्पविराम के स्थान पर~$.$ से अलग करते हैं। यदि $\sigma$ कोई उपसर्ग है, तो
$\sigma.n$, $\sigma \concat \tuple{n}$ है; उदाहरणार्थ यदि
$\sigma = 1.2.1$, तो $\sigma.3$, $1.2.1.3$ है। अतः, मसलन,
\[
\sFmla{\True}{!A \lif (!B \lif !C)}[1.2]
\]
एक \emph{उपसर्गयुक्त !!{signed formula}} (या संक्षेप में केवल
\emph{उपसर्गयुक्त !!{formula}}) है।

सहज रूप से उपसर्ग किसी प्रतिरूप के उस जगत का नाम देता है जो किसी !!{tableau}
की शाखा के !!{formula}ों को संतुष्ट कर सकता है। यदि $\sigma$ किसी जगत का नाम
है, तो $\sigma.n$,~$\sigma$ द्वारा नामित जगत से अभिगम्य किसी जगत का नाम है।

अंतःप्रज्ञावादी प्रतिरूपों में अभिगम्यता संबंध स्वपरावर्ती और संक्रमणशील
होता है। उपसर्गों के पदों में इसका अर्थ है कि $\sigma$ स्वयं~$\sigma$ से
अभिगम्य है, और~$\sigma$ को विस्तृत करने वाला प्रत्येक उपसर्ग भी; अर्थात
$\sigma.n_1.\cdots.n_k$ रूप वाला कोई भी उपसर्ग। स्वयं~$\sigma$ और उसके किसी
भी विस्तार को दर्शाने के लिए $\sigma.*$ संकेत प्रस्तुत करें। दूसरे शब्दों
में, $\sigma.*$ उपसर्ग ठीक वे सभी उपसर्ग हैं जो~$\sigma$ से अभिगम्य हैं।

\end{document}
